\documentclass[12pt]{article}
\usepackage[utf8]{inputenc}
\usepackage[T1]{fontenc}
\usepackage{hyperref}
\title{Tefaro Language Resources for K-12 Curriculum Development}
\author{Compiled by Claude Sonnet 5 (claude-sonnet-5) for Andr\'as Kornai}
\date{}
\begin{document}
\maketitle
\section{Introduction}
Tefaro (ISO 639-3 tfo; also called Demba) is a poorly documented Papuan language of Waropen Regency, on the eastern shore of Cenderawasih (Geelvink) Bay in Indonesia. Wikipedia gives about 100 speakers (1987); Joshua Project gives about 200 people. The search found no educational materials.
\section{Primary Educational Resources}
None found.
\section{Linguistic and Reference Materials}
\begin{itemize}
\item Overview article: \url{https://en.wikipedia.org/wiki/Tefaro_language}
\item Endangered Languages Project entry (with a bibliography page): \url{https://endangeredlanguages.com/lang/5093?hl=en}
\end{itemize}
\section{Religious and Community Resources}
Joshua Project reports that Bible translation has started: \url{https://joshuaproject.net/index.php/people_groups/15317}
\section{Digital and Audio Resources}
None found.
\section{Curriculum Development Context}
Documentation is minimal, so any curriculum would depend on new fieldwork and community-led description first.
\end{document}